\appendix
\section{Notations and formulas}
\label{aa}
 In this sections, we summarize the notations
 and formulas which we frequently use in this paper.
\subsection{Partitions}
\label{a1}
 A partition $\l$
 is a sequence of nonnegative integers
 $\l=(\l_1,\l_2,\dots)$, $\l_j\in\Z_{\ge 0}$ with non decreasing orders $\l_1\ge\l_2\ge\cdots$. The length
 of partition $\ell(\l)$ is the number of nonzero
 elements in $\l$ and the size of $\l$ is defined
 as $|\l|=\sum_{j}\l_j$. Let $\Pe$ be the set of
 the partitions whose lengths are finite.
 The partition $\l\in\Pe$
 can be expressed in another way
 $\l=1^{m_1}2^{m_2}\cdots$, where~$m_j$ represents the number of
 the elements $\l_k$s such that
 $\l_k=j$. Note that $\sum_jm_j=n$
 and $|\l|=\sum_{j} jm_j$.
 For example, $\l=(3,3,2,2,1,1,1,0,\dots)$ is
 expressed as $1^3 2^2 3^2$. We write $\lambda=\varnothing$ if all $\lambda_j$'s are~0's.
 For $\l,\mu\in\Pe$, we denote $\mu\subset\l$ if
 $\mu_i\le\l_i,~i=1,2,\dots$.

\subsection[q-Series and sum over partitions]{$\boldsymbol{q}$-Series and sum over partitions}
\label{a2}
 Here we summarize a few formulas for $q$-series and related summations over partitions. For more detail, we refer the reader to \cite[Sections~10 and~11]{AndrewsAskeyRoy2000}. Define
\begin{align}
\label{a2-1}
 &(a;q)_\i = \prod_{j=0}^\i \big(1-a q^j\big),
 \\
\label{a2-1d}
 &(a;q)_n = \frac{(a;q)_\i}{(a q^n;q)_\i}
 =
 \begin{cases}
 (1-a)(1-aq)\cdots \big(1-a q^{n-1}\big), & n\ge 0,
 \\
 \dfrac{1}{(1-a/q)(1-a/q^2)\cdots(1-a/q^m)},& n=-m,~m>0.
 \end{cases}
\end{align}
Note that $(q;q)_{-m}^{-1}=0$ for $m\in\Z_{>0}$.


 Next we give some formulas for the $q$-series. First, we state the $q$-binomial theorem:
 For $|z|<1$, $|q|<1$ and $a\in\C$,
\begin{align}
\label{a2-2}
\sum_{n=0}^\i \frac{(a;q)_n}{(q;q)_n} z^n
=
\frac{(az;q)_\i}{(z;q)_{\i}}.
\end{align}
For the proof, see, e.g., \cite[Theorem 10.2.1]{AndrewsAskeyRoy2000}. In this paper,
we use the two special cases,
\begin{align}
\label{a2-3}
&\sum_{n=0}^\i \frac{z^n}{(q;q)_n}
=
\frac{1}{(z;q)_{\i}},
\\
 \label{a2-4}
&\sum_{n=0}^\i \frac{(-1)^n q^{\frac{n(n-1)}{2}}}{(q;q)_n} z^n
= (z;q)_\i.
\end{align}
 The first equality \eqref{a2-3} holds for $|z|<1$ while the second one \eqref{a2-4} holds for arbitrary $z\in\C$. They are written in~\cite[Corollary
 10.2.2\,(a) and (b)]{AndrewsAskeyRoy2000}.
 One immediately obtains \eqref{a2-3} (resp.~\eqref{a2-4})
 by setting $a=0$ (resp.\ replacing
 $a\rightarrow 1/a$ $x\rightarrow ax$, then taking the limit $a\rightarrow 0$) in \eqref{a2-2}.

 The Ramanujan's summation formula is a
 generalization of the $q$-Binomial
 theorem~\eqref{a2-2} to the sum over
 $\Z$: For $|q|<1$, $|b/a|<|z|<1$,
\begin{equation}
\label{a2-5}
\sum_{n\in\Z} \frac{(a;q)_n}{(b;q)_n} z^n
=
\frac{(az;q)_\i \big(\frac{q}{az};q\big)_\i (q;q)_\i \big(\frac{b}{a};q\big)_\i}
 {(z;q)_\i \big(\frac{q}{a};q\big)_\i (b;q)_\i \big(\frac{b}{az};q\big)_\i}.
\end{equation}
 For a proof, see, e.g., in \cite[Theorem 10.5.1]{AndrewsAskeyRoy2000}.
 In this paper, we use the special case
 $a=-t$, $b=-qt$, $z=q/w$ with
 $t\in \C\setminus\{0\}$,
\begin{align}
 \label{a2-6}
 \sum_{n\in\Z}\frac{t q^n w^{-n}}{1+tq^n}
 =
 \frac{\theta(-w/t)(q;q)_{\i}^2}
 {\theta(-1/t)\theta(w)},
\end{align}
 where $\theta (x)$ is defined by using
 \eqref{a2-1} as $\theta(x)=(x;q)_{\i}(q/x;q)_{\i}$
 and
 according to the condition on $z$ in \eqref{a2-5}, $w$ should satisfy $|q|<|w|<1$.

Finally, we present a well-known formula for the generating function of $\lambda$:
\begin{equation}
 \label{a2-8}
 \sum_{\l\in\Pe}
 q^{|\l|}=\frac{1}{(q;q)_\i}
\end{equation}
and its refinement with a fixed $\lambda_1$:
\begin{align}
\label{a2-7}
 &\sum_{\substack{\l\in\Pe\\ \l_1=n}}
 q^{|\l|}=\frac{q^n}{(q;q)_n}.
% \\
\end{align}
 \eqref{a2-7} is obtained as follows. Switching to the notation of partition as $\l=1^{m_1}2^{m_2}\cdots$,
 and noting $|\l|=\sum_{j}jm_j$, we have
\begin{align*}
 \sum_{\substack{\l\in\Pe\\ \l_1=n}}
 q^{|\l|}
 =
 \prod_{j=1}^{n-1}\sum_{m_j=0}^{\i}
 q^{jm_j}\cdot \sum_{m_n=1}^{\i}q^{n m_n}=\prod_{j=1}^{n-1}\frac{1}{1-q^j}
 \cdot\frac{q^n}{1-q^n}.
\end{align*}
\eqref{a2-8} follows immediately from \eqref{a2-7} by taking the sum over $n$ and noting \eqref{a2-3}.



\section[Properties of the kernels A(m;r) and B(r;m)]{Properties of the kernels $\boldsymbol{A(m;r)}$ and $\boldsymbol{B(r;m)}$}
\label{ab}
 In this appendix, we give upper bounds for matrix elements of $A(m;r)$ \eqref{3-15d} and $B(r;m)$ \eqref{3-16d}, which are useful to prove the well-definedness of the Fredholm determinants $F_{\ell}$\eqref{d31-1} and $F_{\i}$\eqref{d31-2} in Proposition~\ref{p33}\,(ii) and Lemma~\ref{l32}\,(ii) in Section~\ref{s33}.
\begin{Lemma}
\label{lb1}
Recall the assumption on $a_i$, $b_i$, $i=1,\dots,N$ stated at the beginning of Section~{\rm\ref{s3}}. Recall also the notation \eqref{s32-1}, $\tilde{a}_r=a_k q^{u}$. Then there exist constants $C_+$, $C_-$, $D_+$, and $D_-$ which do not depend on $m$ or
$r$ such that the following inequalities hold.
\begin{itemize}\itemsep=0pt
 \item [$(i)$] Assume $b$ satisfies $b_{\max}<b<1$. We have
 \begin{align}
 \label{bb2-1}
 |A(m;r)|
 \begin{cases}
 < C_+ a_{1}^mq^{\frac{m^2}{2N}+\frac{m}{2}+u}, & m\ge 0,
 \\
 < C_- b^{-m} q^{u}, & m<0.
 \end{cases}
 \end{align}
 \item[$(ii)$] Assume $a_{1}/a_{N}<q^{-\frac12+\e}$ with $\e\in(0,1/2)$. We have
 \begin{align}
 \label{bb2-2}
 |B(r;m)|
 \begin{cases}
 < D_+ a_{N}^{-m}q^{-\frac{m^2}{2N}+\e m-(\frac12+\e)u}, & m\ge 0,
 \\
 < D_- a_{1}^{-m}q^{u}, & m<0.
 \end{cases}
 \end{align}
\end{itemize}
\end{Lemma}

From \eqref{bb2-1}, \eqref{bb2-2}, we immediately obtain the bounds
for $\tilde{A}(m;r)$ and $\tilde{B}(r;m)$ defined by~\eqref{3-16dd}.

\begin{Corollary}\label{cb2}
With the same notation and assumptions as in Lemma~{\rm\ref{lb1}},
we have
\begin{align*}%\label{bb2-2d}
 \big|\tilde{A}(m;r)\big|
 \begin{cases}
 < C_+ q^{\e m+\o u}, & m\ge 0,
 \\
 < C_- b^{-m} q^{\omega u}, & m<0,
 \end{cases}
 \qquad
 \big|\tilde{B}(r;m)\big|
 \begin{cases}
 < D_+ q^{\e m+(\frac12-\e-\o)u}, & m\ge 0,
 \\
 < D_- a_1^{-m}q^{u}, & m<0.
 \end{cases}
 \end{align*}
 \end{Corollary}

 Note that
 both $\tilde{A}(m;r)$ and $\tilde{B}(r;m)$ are bounded exponentially for $m$ and $r$.

\begin{proof}[Proof of Lemma~\ref{lb1}]
In this proof, $C_i$, $i=1,\dots,4$ and $D_j$, $j=1,2,3$ will denote constants which do not depend on $m$ or $r$ (i.e., $k$, $u$).

(i) First, we consider the case $m\ge 0$.
We note that the singularities of the integrand of~\eqref{3-15d} are~$0$ and $b_iq^j,$ $i=1,\dots,N$, $j\in\Z_{\ge0}$ and do not include $\tilde{a}_r^{-1}$
because of the presence of the factor $\prod_{i=1}^N(a_iz;q)_{\i}$. Thus we can extend the contour $C$ arbitrarily without changing
the value of $A(m;r)$. First we give an estimate setting the radius to be $q^{-n}$ with $n\in\Z_{>0}$ and then specify~$n$ giving an optimal estimate. With this choice of the radius $|z|=q^{-n}$, $f(m)$ in \eqref{p251-1} and some factors in the integrand in $A(m;r)$
can be estimated as
\begin{gather}
 |f(m)|<q^m,
\qquad
 \frac{2\pi q^{-n}}{|z^{1+m}|}=2\pi q^{n m},\nonumber
\\
 \left|\prod_{j=1}^N\frac{1}{(b_j/z;q)_{\i}}\right|
 <\prod_{j=1}^N\frac{1}{(b_j q^n;q)_{\i}}
 <\prod_{j=1}^N\frac{1}{(b_j;q)_{\i}},\label{bb2-3}
\end{gather}
where the factor $2\pi q^{-n}$ represents
the circumference with radius $q^{-n}$.
We rewrite the remaining factor as
\begin{align}
\label{bb2-4}
 \frac{1}{z-\tilde{a}_{r}^{-1}}
 \prod_{i=1}^N(a_i z;q)_{\i}
 =
 -a_k q^{u}
 \prod_{\substack{i=1\\i\neq k}}^N(a_i z;q)_{\i}
 \cdot (a_k z;q)_{u} \big(a_k q^{u+1}z;q\big)_{\i}.
\end{align}
The factor \smash{$-a_kq^u\prod_{\substack{i=1\\i\neq k}}^N(a_i z;q)_{\i}$} in \eqref{bb2-4} can be estimated as
\begin{gather}
 \Biggl|-a_kq^u\prod_{\substack{i=1\\i\neq k}}^N(a_i z;q)_{\i}\Biggr|\nonumber\\
\qquad < a_kq^u\prod_{\substack{i=1\\i\neq k}}^N\bigl(-a_i/q^n ;q\bigr)_{\i}
 =a_k
 q^{u-(N-1)\frac{n(n+1)}{2}}
 \prod_{\substack{i=1\\i\neq k}}^N a_i^n
 (-q/a_i;q)_n(-a_i;q)_\i
 \nonumber\\
\qquad <a_k
 q^{u-(N-1)\frac{n(n+1)}{2}}
 \prod_{\substack{i=1\\i\neq k}}^N a_i^n
 (-q/a_i;q)_\infty (-a_i;q)_\i
 <C_1 q^{u-(N-1)\frac{n(n+1)}{2}} a_{1}^{(N-1)n},\label{bb2-5}
\end{gather}
 where
 in the second equality we used the relation
 \begin{align}
 \label{bb2-6}
 \big(x/q^n;q\big)_\i=(-x)^n q^{-\frac{n(n+1)}{2}}
 (q/x;q)_{n}(x;q)_{\i}.
 \end{align}
 For the remaining factors in \eqref{bb2-4}, we find
 \begin{gather*}% \label{bb2-7}
 \big|(a_k z;q)_u \big(a_k q^{u+1}z;q\big)_{\i}\big| \\
 \qquad <\bigl(-a_k /q^n;q\bigr)_u\bigl(-a_k q^{u-n+1};q\bigr)_{\i}=\frac{\bigl(-a_k q^{-n};q\bigr)_\infty}{1+a_k q^{u-n}} =
 a_k^n q^{-\frac{n(n+1)}{2}}\!\frac{(-q/z_k;q)_n(-a_k;q)_{\infty}}{1+a_kq^{u-n}}\\
 \qquad <
 \begin{cases}
 C_2 a_k^n q^{-\frac{n(n+1)}{2}}, & u\ge n,
 \\
 C_3 a_k^n q^{-\frac{n(n+1)}{2}+n-u}, & n>u.
 \end{cases}
 \end{gather*}
 Here in the third equality we used \eqref{bb2-6}.
 Thus we have a bound regardless of the order of $n$ and~$u$
 \begin{align}
 \label{bb2-8}
 \big|(a_k z;q)_u\big(a_kq^{u+1}z;q\big)_{\i}\big|
 <
 C_4 a_{1}^nq^{-\frac{n(n+1)}{2}}.
 \end{align}
 Putting together \eqref{bb2-3}, \eqref{bb2-5}, and \eqref{bb2-8}, we have
 \begin{align}
 \label{bb2-9}
 |A(m;r)|&<C_+ a_{1}^{Nn} q^{nm-N\frac{n(n+1)}{2}+u+m}
 \nonumber\\
 &=C_+ a_{1}^{Nn}q^{-\frac{N}{2}\left(n-\frac{m}{N}+\frac12\right)^2+\frac{N}{2}\left(\frac{m}{N}-\frac12\right)^2+u+m}.
 \end{align}

 Now we specify $n$ which gives an optimal estimate.
 From the last expression in \eqref{bb2-9}, a~simple computation shows that such optimal bound is given by the nonnegative integer which is closest to $m/N-1/2$ and this proves the first inequality in \eqref{bb2-1}.


 Next we consider the case $m<0$. In this case, we set the radius of contour $C$ to be $b$ appearing in the assumption of (i). We find{\samepage
 \begin{align*}
 &|f(m)|<C_4
 \qquad
 \frac{2\pi b}{|z|^{1+m}}=2\pi b^{-m},
 \qquad
 \frac{1}{\big|z-\tilde{a}^{-1}_r\big|}
 <\frac{a_k q^{u}}{1-a_k bq^{u}}
 <\frac{a_k}{1-a_k} q^{\ell},
 \nonumber\\
 &\Bigg|\prod_{j=1}^M\frac{(a_iz;q)_{\i}}{(b_j/z;q)_{\i}}\Bigg|
 <\prod_{i=1}^N\frac{(-a_ib;q)_\i}{(b_i/b;q)_{\i}}.
 \end{align*}
 Combining these estimates, we arrive at the bottom estimate of \eqref{bb2-1}.}


 (ii)
 Recalling the notation $\tilde{a}_r=a_k q^{u}$ \eqref{s32-1} and calculating explicitly the residue in \eqref{3-16d}, we have
 \begin{align*}
 B(r;m)=&-a_k^{-(m+1)}q^{-u(m+1)}
 \prod_{\substack{i=1\\ i\neq k}}
 \frac{1}{(a_iq^{-u}/a_k ;q)_{\i}}\cdot
 \frac{1}{(1/q^u;q)_u (q;q)_{\i}}\cdot
 \prod_{j=1}^N(a_kb_jq^{u};q)_{\i}
 \\
 =&
 -a_k^{-(m+1)}q^{-u (m+1)+N\frac{u(u+1)}{2}}
 \prod_{\substack{i=1\\ i\neq k}}
 \frac{(-a_k/a_i)^u}{(q a_k/a_i;q)_{\i}(a_i/a_k;q)_\infty}\cdot
 \frac{(-1)^u}{(q;q)_u (q;q)_{\i}}
 \\
 &\times\prod_{j=1}^N(a_k b_jq^{u};q)_{\i},
 \end{align*}
 where in the second equality we used \eqref{bb2-6}.
 Using the condition $a_1/a_N< q^{-1/2+\e}$, we have the following estimate,
 \begin{align}
 \label{bb2-12}
 |B(r;m)|&<D_1 a_k^{-(m+1)} q^{-u m +N\frac{u(u+1)}{2}-N\left(\frac12-\e\right)u -\left(\frac12+\e\right)u}.
 \end{align}
 For $m\ge 0$, combining \eqref{bb2-12} with
 the following bound
 \begin{align*}
 q^{-u m+N\frac{u(u+1)}{2}-N\left(\frac12-\e\right)u}
 =
 q^{\frac{N}{2}\left(u-\frac{m}{N}+\e\right)^2-\frac{N}{2}\left(\frac{m}{N}-\e\right)^2}
 <D_2
 q^{-\frac{m^2}{2N}+\e m},
 \end{align*}
 we have the desired estimate. On the other hand
 for $m<0$, we use
 \begin{align}
 \label{bb2-14}
 q^{-u m}<q^{u},
 \qquad
 q^{N\frac{u(u+1)}{2}-N\left(\frac12-\e\right)u -\left(\frac12+\e\right)u}<D_3.
 \end{align}
 From \eqref{bb2-12} and \eqref{bb2-14}, we obtain
 the bottom relation in \eqref{bb2-2}.
\end{proof}



\section[Decay estimates of q-Pochhammer symbols]{Decay estimates of $\boldsymbol{q}$-Pochhammer symbols}
\label{ac}
In this section, we provide simple bounds for
the $q$-Pochhammer symbol $(z;q)_{\i}$, where
$z\in\C$ and $0<q<1$. For our purpose, it is sufficient
to have bounds in the following three cases though
they are not exclusive nor exhaustive.
\begin{Lemma}\label{lc1}
 Let us write $z=a+\sqrt{-1}b$ with $a,b\in\R$.
 For $0<q<1$ fixed, there exist positive constants
 $c_1=c_1(q)$ and $c_2=c_2(q)$ such that
 \begin{itemize}\itemsep=0pt
 \item[\rm{(i)}] for $a\le -1$,
 \begin{align}
 \label{lc1-1}
 |(z;q)_\i|\ge c_1 \exp \big[c_2\log^2 |a|\big],
 \end{align}
 \item[\rm{(ii)}] for $a<1$ and $|b|>1$,
 \begin{align}
 \label{lc1-2}
 |(z;q)_\i|\ge c_1 \exp\big[c_2\log^2 |b|\big],
 \end{align}
 \item[\rm{(iii)}] for $a>1$,
 the following two types of lower bounds
 hold
 \begin{align}
 \label{lc1-3}
 |(z;q)_\i|
 &\ge c_1 b^2 a^{-\frac32}
 \exp\big[c_2\log^2 a\big],\\
 \label{lc1-3'}
 |(z;q)_\i|
 &\ge c_1 \big(q^{\a-1}-1\big)\big(1-q^{\a}\big)
 \exp\big[c_2\log^2 a\big],
 \end{align}
 where $\a=\left\lceil \log_{q^{-1}}a\right\rceil-\log_{q^{-1}}a$
 with $\lceil x\rceil$ being the minimum integer
 greater than $x$.
 \end{itemize}
\end{Lemma}


\begin{proof}
 We repeatedly use the simple estimate
 \begin{align}
\label{lc1-5}
 |1-w|=\sqrt{(1-\Re w)^2+(\Im w)^2}\ge
 \begin{cases}
 |1-\Re w|,\\
 |\Im w|,
 \end{cases}
\end{align}
 which hold for any $w\in\C$. We use the top or bottom estimate depending on the situations.

 (i) In terms of $\a\in [0,1)$ and $J\in\Z_{>0}$, we express $a(<-1)$
 as $a=-q^{\a-J}$.
 Applying the top estimate in \eqref{lc1-5}
 to the factor $\big|1-z q^k\big|$, we have
 \begin{align}
 \label{lc1-6}
 \big|1-z q^k\big|=\sqrt{\big(1+aq^{k}\big)^2+\big(bq^k\big)^2}
 \ge 1+q^{\a-J+k}
 \ge
 \begin{cases}
 q^{\a-J+k},
 \\
 1-q^{k+1-J},
 \end{cases}
 \end{align}
 where both estimates holds for $\forall k=0,1,2,\dots$. Choosing the top (resp.\ bottom) estimate for~${k<J}$ (resp.\ $k\ge J$), we have
 \begin{align}
 \label{lc1-7}
 |(z;q)_\i|\ge q^{\a J-\frac{J(J+1)}{2}}(q;q)_{\i}
 =
 q^{-\frac{1}{2}(\a-J)^2+\frac{1}{2}(\a-J)
 +\frac{\a^2-\a}{2}}(q;q)_{\i}
 \ge
 q^{-\frac{1}{2}(\a-J)^2}(q;q)_{\i},
 \end{align}
 where in the last inequality we used the fact
 $0\le\a<1$ and $1\le J$. Recalling $\a-J=\log_q|a|$,
 one easily sees
 \begin{align}
 \label{lc1-8}
 q^{-\frac{(\a-J)^2}{2}}=
 {\rm e}^{\frac{-(\log |a|)^2}{2\log q}}.
 \end{align}
 From \eqref{lc1-7} and \eqref{lc1-8},
 we obtain \eqref{lc1-1}.

 (ii)
 We adopt the same strategy as above.
 In this case we write $|b|=q^{\a-J}$, where
 $0\le \a<1$ and $J\in\Z_{>0}$. Using \eqref{lc1-5},
 we have for $k=0,1,2,\dots$
 \begin{align}
 \label{lc1-9}
 \big|1-z q^k\big|
 \ge
 \begin{cases}
 |b|q^{k}=q^{\a-J+k},
 \\
 1-aq^k\ge 1-q^{k+1-J}.
 \end{cases}
 \end{align}
 Applying the top (resp.\ bottom) estimate in \eqref{lc1-9} to the case $k<J$ (resp.\ $k\ge J$) and noting that the estimate \eqref{lc1-9}
 is exactly the same as \eqref{lc1-6} except
 that $\a-J$ represents $\log_q|b|$, we get~\eqref{lc1-2} similarly to the case
 (i).

 (iii)
 As in the case (i),
 we express $a$ as $a=q^{\a-J}$ with $a\in[0,1)$ and $J\in\{1,2,\dots\}$.
 First, we prove \eqref{lc1-3}.
 We choose the estimates as follows:
 \begin{gather}
 \big|1-z q^k\big|\nonumber
 \\
 \qquad\ge
 \begin{cases}
 \big|1-a q^k\big|=q^{\a-J+k}\big(1-q^{J-\a-k}\big)&\\
 \phantom{ \big|1-a q^k\big|}{} \ge q^{\a-J+k}\big(1-q^{J-k-1}\big), & {\text{for $0\le k\le J-2$}},
 \\
 bq^k, & {\text{for $k=J-1,J$}},
 \\
 \big|1-a q^k\big|=1-q^{\a-J+k}\ge 1-q^{-J+k}, & {\text{for $J+1\le k$}}.
 \end{cases}\label{lc1-10}
 \end{gather}
 Thus we have
 \begin{align*}% \label{lc1-11}
 |(z;q)_{\i}|
 &\ge
 \prod_{k=0}^{J-2} q^{\a-J+k}\big(1-q^{J-k-1}\big)
 \cdot
 \prod_{k=J-1}^{J} bq^{k}
 \cdot
 \prod_{k=J+1}^{\i} \big(1-q^{-J+k}\big)
 \\
 &=
 q^{\a(J-1)-\frac{J(J+1)}{2}+2J}b^2(q;q)_{J-1}(q;q)_{\infty}\ge
 q^{-\frac{(J-\alpha)^2}{2}+\frac32(J-\alpha)+\frac{\alpha^2}{4}+\frac{3\alpha}{2}}b^2(q;q)_{\infty}^2.
 \end{align*}
 Thus from $a=q^{\alpha-J}$ and \eqref{lc1-8}, we
 obtain the desired form \eqref{lc1-3}.

 The other estimate \eqref{lc1-3'}
 can be readily
 obtained by replacing the second inequality in \eqref{lc1-10}
 (corresponding to the case $k=J-1, J$)
 with $\big|1-zq^k\big|\ge \big|1-aq^k\big|=\big|1-q^{\a-J+k}\big|$.
\end{proof}

